% SGA 3, Expose III, Loop 1 English reconstruction.
% Scope: complete section 3, from its heading through Corollary 3.10,
% stopping immediately before section 4.
% Authority: Polo--Gille Exp3-14oct24.pdf, component-local pp. 37--42
% (combined-reader pp. 147--152; original printed pp. 121--126), with
% only the section-3 portions of the first and last pages admitted.
% Force-OCR is a locator only. English wording was compared with Jacob C.
% Reinhold's CC BY 4.0 Markdown lineage at commit
% e7a259f3f8608ad3edf9bf6eead3fd504dd2d23e.
% The Polo--Gille PDF controls prose, mathematics, numbering, and notes.

\setcounter{footnote}{64}

% Source page: Exp3 local p. 37; combined p. 147; original printed p. 121.
\section*{\SGAInlineRefTarget{sga3:iii:section:3}3. 群概形的无穷小扩张}
\addcontentsline{toc}{section}{3. 群概形的无穷小扩张}

仍沿用 \S0 的记号 (\(S,\mathscr I,\mathscr J\) 等),
设 \(X\) 是一个 \(S\)-概形, 并假设 \(X_{\mathscr J}\)
配备了一个群结构. 我们要确定 \(X\) 上诱导 \(X_{\mathscr J}\)
上给定结构的 \(S\)-群结构.

从现在起, 假设 \(X\) 在 \(S\) 上平坦. 设 \(\mathscr C\) 是
在 \(S\) 上平坦的 \(S\)-概形所构成的范畴;
于是 \(X\in\operatorname{Ob}\mathscr C\).
以 \(Y\) (分别地, \(M\)) 表示 \(\mathscr C\) 上由
\(X^+\) (分别地, \(L'_X\)) 定义的函子.
典范态射 \(p_X:X\to X^+\) 定义了 \(\widehat{\mathscr C}\) 中的态射
\[
 p:X\longrightarrow Y,
\]
而 \(\widehat{(\mathrm{Sch})/S}\) 中 \(L'_X\) 在 \(X\) 上的作用,
定义了 \(\widehat{\mathscr C}\) 中 \(M\) 在 \(X\) 上的作用.
立即可验证, \(X\) 由此成为 \(Y\) 上关于 \(M_Y\) 的形式主齐性对象
(参见 \SGARefLink{sga3:iii:proposition:0-2:item:i}{0.2(i)} 和
\SGARefLink{sga3:iii:remark:0-4}{0.4}).

% Source page: Exp3 local p. 38; combined p. 148; original printed p. 122.
\SGARefLink{sga3:iii:proposition:0-8}{0.8} 中定义的 \(X^+\) 在
\(L'_X\) 上的作用 (彼处记作 \(\operatorname{Ad}\)),
定义了 \(Y\) 在 \(M\) 上的一个作用, 记作 \(\widetilde f\).
另一方面, 由 \SGARefLink{sga3:iii:remark:0-5}{0.5},
\[
 \operatorname{Hom}_{\widehat{\mathscr C}}(Z,M)
 \simeq \operatorname{Hom}_{S_0}(Z_0,L_0),
 \qquad Z\in\operatorname{Ob}\mathscr C,
\]
其中 \(L_0\) 是 \SGARefLink{sga3:iii:remark:0-5}{0.5} 中定义的函子.

\medskip
\noindent\SGAInlineRefTarget{sga3:iii:lemma:3-1}\textbf{引理 3.1.}
\begin{enumerate}
\item[(i)] \textit{\SGARefLink{sga3:iii:subsection:1-3}{1.3} 的条件
\((+)_{n}\) 对每个正整数 \(n\) 都成立.}
\item[\SGAInlineRefTarget{sga3:iii:lemma:3-1:item:ii}(ii)]
\textit{令 \(S_0\)-群 \(X_0\) 通过其伴随表示作用于
\(S_0\)-函子 \(L_0\). 则有典范同构}
\[
 H^*(X_0,L_0)\simeq H^*(Y,M),
\]
\textit{其中第一个上同调在 \((\mathrm{Sch})/S_0\) 中计算,
第二个上同调在 \(\mathscr C\) 中计算.}
\end{enumerate}

引理的两部分都由关系式
\begin{align*}
 \operatorname{Hom}_{\widehat{\mathscr C}}(Y,M)
 &\simeq
 \operatorname{Hom}_{\widehat{(\mathrm{Sch})/S_0}}
       (X^+\times_S S_0,L_0)\\
 &\simeq \operatorname{Hom}_{S_0}(X_0,L_0)\\
 &\simeq \operatorname{Hom}_{\widehat{\mathscr C}}(X,M),
\end{align*}
得到, 而该关系立即由 \(M\) 作为 \(\prod_{S_0/S}L_0\) 的定义得到.
更一般地, 用 \(X^n,Y^n\) 代替 \(X,Y\) 后同样的关系仍成立.
因此, 每个态射 \(X^n\to M\) 都唯一地通过 \(Y^n\) 分解,
这就给出 \((+)_{n}\). 它还给出关系式
\(C^*(Y,M)=C^*(X_0,L_0)\), 从而得到
\SGARefLink{sga3:iii:lemma:3-1:item:ii}{(ii)}.

因此可以应用 \SGARefLink{sga3:iii:subsection:1-3}{1.3} 中的构造.
特别地:

\medskip
\noindent\SGAInlineRefTarget{sga3:iii:lemma:3-2}\textbf{引理 3.2.}
\textit{设 \(P:X\times_S X\to X\) 是一个态射.
则 \(P\) 诱导 \(X_{\mathscr J}\) 的群法则,
当且仅当 \(P\) 是 \(X\) 上的一个容许合成法则
(参见 \SGARefLink{sga3:iii:definition:1-3-2}{1.3.2}).}

事实上, \(P\) 诱导 \(X_{\mathscr J}\) 的群法则,
当且仅当它提升了 \(X^+\) 的群法则, 或等价地提升了 \(Y\) 的群法则.
所以只需证明, 每个提升 \(X_{\mathscr J}\) 的群法则的态射 \(P\)
都满足 \SGARefLink{sga3:iii:definition:1-3-2:item:ii}{1.3.2(ii)} 的恒等式 ((++)),
而这恰好就是 \SGARefLink{sga3:iii:proposition:0-8}{0.8} 中证明的内容.

\medskip
\noindent\SGAInlineRefTarget{sga3:iii:proposition:3-3}\textbf{命题 3.3.}
\textit{设 \(S\) 是一个概形, \(S_0\) 是由幂零理想定义的闭子概形.
设 \(X\) 是一个在 \(S\) 上平坦, 且在 \(S\) 上拟紧或局部有限呈示的
\(S\)-概形. 设 \(P:X\times_S X\to X\) 是 \(X\) 上的合成法则.
则 \(P\) 是群法则, 当且仅当满足以下两个条件:}
\begin{enumerate}
\item[\SGAInlineRefTarget{sga3:iii:proposition:3-3:item:i}(i)] \textit{\(P\) 是结合的;}
\item[\SGAInlineRefTarget{sga3:iii:proposition:3-3:item:ii}(ii)]
\textit{\(P\) 在 \(X_0=X\times_S S_0\) 上诱导一个群法则.}
\end{enumerate}

这些条件显然是必要的. 下面证明其充分性.
先假设 \(X\to S\) 有一个截面.
此时, 对每个 \(S'\to S\), \(X(S')\) 都非空, 所以只需
\footnote{编者注: 原文已作修正, 删去了对 Bourbaki 关于半群的一道习题的不恰当引用
(参见 [BAlg], \S I.2, 习题 9--13), 并说明了左、右平移的作用;
见下一条编者注.}
证明, 对每个 \(x\in X(S')\), 由 \(x\) 给出的左、右平移都是
\(X_{S'}\) 的同构.
\footnote{编者注: 设 \(E\) 是一个带有结合合成法则的非空集合,
并且每个左平移 \(\ell_x\) 都是双射. 固定 \(x_0\in E\).
存在唯一的 \(e\in E\) 使得 \(x_0e=x_0\);
于是 \(x_0ex=x_0x\) 蕴含对每个 \(x\in E\) 都有 \(ex=x\).
此外, 对每个 \(x\), 存在唯一的 \(x'\) 使得 \(xx'=e\).
再假设存在一个 \(b\in E\), 使得右平移 \(r_b\) 是单射.
则由 \(xeb=xb\) 得到 \(xe=x\), 所以 \(e\) 是单位元;
由 \(xx'x=x=xe\) 得到 \(x'x=e\). 因此 \(x'\) 同时为 \(x\)
的左逆元和右逆元, 而 \(E\) 是一个群.

\(r_b\) 为单射的假设是必要的: 在任意集合 \(E\) 上,
对所有 \(x,y\in E\) 定义 \(xy=y\).
此时每个左平移都是恒等映射 (且该法则是结合的),
而对每个 \(y\), 有 \(r_y(E)=\{y\}\);
所以当 \(|E|>1\) 时, \(E\) 不是群.}

% Source page: Exp3 local p. 39; combined p. 149; original printed pp. 123--124.
显然可以假设 \(S'=S\). 所讨论的平移 \(t\) 在 \(X_0\) 上诱导
一个平移 \(t_0\), 由于 \(X_0\) 是群, 它是一个自同构.
由平坦性即得结论 (SGA 1 III, 4.2).
\footnote{编者注: 由于 \(X\) 与 \(X_0\) 有相同的底拓扑空间,
且 \(t_0\) 是自同构, 所以 \(t\) 是同胚, 从而是仿射态射;
参见第 VI\textsubscript{B} 篇, 2.9.1, 或 EGA IV\textsubscript{4}, 18.12.7.1.
因此只需证明如下断言. 设 \(\mathscr J\) 是环 \(\Lambda\) 的幂零理想,
\(\varphi:A\to B\) 是一个 \(\Lambda\)-代数态射,
其中 \(B\) 在 \(\Lambda\) 上平坦,
且 \(\varphi\otimes_\Lambda(\Lambda/\mathscr J)\) 是双射.
则 \(\varphi\) 是双射. 由 ``幂零 Nakayama 引理'',
\(\varphi\) 是满射; 此外, 由于 \(B\) 在 \(\Lambda\) 上平坦,
还有 \(\operatorname{Ker}(\varphi)\otimes_\Lambda
(\Lambda/\mathscr J)=0\), 从而
\(\operatorname{Ker}(\varphi)=0\). 因此 \(\varphi\) 是双射.}

现在去掉 \(X\) 在 \(S\) 上有截面的假设,
并假设存在 \(S'\to S\), 使得 \(X_{S'}\) 在 \(S'\) 上有一个截面.
由刚刚证明的结果, \(X_{S'}\) 是一个 \(S'\)-群;
设 \(e'\) 是其单位截面.
\(e'\) 在 \(\operatorname{pr}_i:S'\times_S S'\to S'\), \(i=1,2\)
下的逆像, 是 \(X_{S''}\) 关于将 \(P_{S'}\) 沿
\(\operatorname{pr}_i\) 拉回所得的群法则的单位截面.
然而, 由于 \(P\) ``定义在 \(S\) 上'', 这两个群法则一致,
其单位截面也一致. 因而
\[
 \operatorname{pr}_1^*(e')=\operatorname{pr}_2^*(e').
\]

若 \(S'\to S\) 是下降态射 (参见第 IV 篇, 第 2 号),
则存在 \(X\) 的一个截面, 其基扩张为 \(e'\), 证明即告完成.
由于 \(X_X\) 在 \(X\) 上有一个截面, 即对角截面,
现在只需证明 \(X\to S\) 是下降态射.
它是平坦满射, 且拟紧或局部有限呈示;
因此它是一个 (fpqc) 覆盖, 从而是下降态射 (第 IV 篇, 第 6 号).

\medskip
\noindent\textbf{注.}
事实上, 由以下结果可知, \(X\to S\) 拟紧或局部有限呈示的假设是多余的.
读者可将这一结果作为第 IV 篇的一道习题来证明:

在对 \(S\) 和 \(S_0\) 所作的上述假设下,
若 \(X\to S\) 平坦且 \(X_0\to S_0\) 是一个 (fpqc) 覆盖,
则 \(X\to S\) 是下降态射.

\medskip
\noindent\SGAInlineRefTarget{sga3:iii:lemma:3-4}\textbf{引理 3.4.}
\textit{\(X\) 上的两个容许合成法则等价
(参见 \SGARefLink{sga3:iii:definition:1-3-5}{1.3.5}),
当且仅当其中一个由另一个通过 \(X\) 的某个 \(S\)-自同构得到,
该自同构在 \(X_{\mathscr J}\) 上诱导恒等映射.}

% Source page: Exp3 local p. 40; combined p. 150; original printed p. 125.
事实上, 由 \SGARefLink{sga3:iii:corollary:0-7}{0.7},
\SGARefLink{sga3:iii:lemma:1-3-1}{1.3.1} 中构造的态射恰好就是上述断言中的态射.
\footnote{编者注: 事实上, 由 \SGARefLink{sga3:iii:corollary:0-7}{0.7} 的证明,
\(X\) 上在 \(X_{\mathscr J}\) 上诱导恒等映射的 \(S\)-自同态,
是形如 \(m\cdot\operatorname{id}_X\) 的自同构,
其中 \(m\) 遍历 \(M(X)=\operatorname{Hom}_S(X,M)\).
对每个 \(S'\to S\) 及每个 \(x\in X(S')\), 有
\((m\cdot\operatorname{id}_X)(x)=m(x)\cdot x\).
由 \SGARefLink{sga3:iii:lemma:3-1}{3.1} 的证明,
每个 \(m:X\to M\) 都唯一地通过某个态射 \(h:Y=X^+\to M\) 分解;
所以 \(m\cdot\operatorname{id}_X\) 就是
\SGARefLink{sga3:iii:lemma:1-3-1}{1.3.1} 中引入的自同构 \(u_h\).
由等价关系的定义便得此引理;
参见 \SGARefLink{sga3:iii:lemma:1-3-4}{1.3.4} 和
\SGARefLink{sga3:iii:definition:1-3-5}{1.3.5}.}

综合以上所有结果, 由\SGARefLink{sga3:iii:proposition:1-3-6}{命题 1.3.6} 得到:

\medskip
\noindent\SGAInlineRefTarget{sga3:iii:theorem:3-5}\textbf{定理 3.5.}
\textit{设 \(S\) 是一个概形, \(\mathscr I\) 与 \(\mathscr J\)
是 \(S\) 上的两个理想, 使得 \(\mathscr I\supset\mathscr J\)
且 \(\mathscr I\mathscr J=0\).
设 \(S_0\) 与 \(S_{\mathscr J}\) 是它们定义的闭子概形.
设 \(X\) 是一个在 \(S\) 上平坦 (且在 \(S\) 上局部有限呈示或拟紧)
的 \(S\)-概形, 以 \(X_0\) 和 \(X_{\mathscr J}\) 表示基变换所得的概形.
假设 \(X_{\mathscr J}\) 配备了一个 \(S_{\mathscr J}\)-群结构,
以 \(L_0\) 表示由下式定义的取值于交换群的 \(S_0\)-函子:}
\[
 \operatorname{Hom}_{S_0}(T,L_0)
 =\operatorname{Hom}_{\mathscr O_T}
 \left(
  \omega^1_{X_0/S_0}\otimes_{\mathscr O_{S_0}}\mathscr O_T,
  \mathscr J\otimes_{\mathscr O_{S_0}}\mathscr O_T
 \right),
\]
\textit{\(X_0\) 通过其伴随表示作用于该函子.}
\begin{enumerate}
\item[\SGAInlineRefTarget{sga3:iii:theorem:3-5:item:i}(i)]
\textit{\(X\) 上存在诱导 \(X_{\mathscr J}\) 上给定结构的
\(S\)-群结构, 当且仅当满足以下条件:}
  \begin{enumerate}
  \item[\((i_1)\)] \textit{存在一个 \(S\)-概形态射
  \(X\times_S X\to X\), 诱导 \(X_{\mathscr J}\) 的群法则;}
  \item[\((i_2)\)] \textit{由 \(X\) 及 \(X_{\mathscr J}\) 的群法则
  典范定义的 \(H^3(X_0,L_0)\) 中的某个障碍类为零.}
  \end{enumerate}
\item[\SGAInlineRefTarget{sga3:iii:theorem:3-5:item:ii}(ii)]
\textit{若 \SGARefLink{sga3:iii:theorem:3-5:item:i}{(i)} 中的条件成立,
则 \(X\) 上诱导 \(X_{\mathscr J}\) 上给定法则的群法则所成的集合
\(\mathcal E\), 是 \(Z^2(X_0,L_0)\) 下的主齐性集合.
此外, \(\mathcal E\) 关于 \(X\) 的那些在 \(X_{\mathscr J}\)
上诱导恒等映射的 \(S\)-自同构的商,
是 \(H^2(X_0,L_0)\) 下的主齐性集合.}
\end{enumerate}

\footnote{编者注: 以下论证为新增内容.}
事实上, 由 \SGARefLink{sga3:iii:lemma:3-2}{3.2}, 每个诱导
\(X_{\mathscr J}\) 的群法则的 \(S\)-概形态射
\(f:X\times_S X\to X\) 都是 \(X\) 上的容许合成法则.
由 \SGARefLink{sga3:iii:proposition:1-3-6:item:i}{1.3.6(i)},
存在一个结合的容许合成法则 \(P:X\times_S X\to X\),
等价于某个类 \(c(f)\in H^3(X_0,L_0)\) 消失;
此时, \SGARefLink{sga3:iii:proposition:3-3}{3.3} 表明 \(P\) 是群法则.
这证明了 \SGARefLink{sga3:iii:theorem:3-5:item:i}{(i)},
而 \SGARefLink{sga3:iii:theorem:3-5:item:ii}{(ii)} 由
\SGARefLink{sga3:iii:proposition:3-3}{3.3} 及
\SGARefLink{sga3:iii:proposition:1-3-6:item:ii}{1.3.6(ii)} 得到.

\medskip
\noindent\SGAInlineRefTarget{sga3:iii:remark:3-5-1}\textbf{注 3.5.1.}
\footnote{编者注: 此注与 \SGARefLink{sga3:iii:remark:4-5-1}{4.5.1} 类似,
为新增内容, 用以引入 \SGARefLink{sga3:iii:paragraph:4-38}{4.38}
中使用的记号 \(\delta(\mu,\mu')\) 或 \(\delta(X,X')\).
因此, \SGARefLink{sga3:iii:theorem:3-5:item:ii}{3.5(ii)} 中
关于 \(\mathcal E\) 本身的部分也是新增的.}
若 \(\mu\) 与 \(\mu'\) 是 \(X\) 上诱导 \(X_{\mathscr J}\) 上给定法则的群法则,
则得到一个余循环 \(\delta(\mu,\mu')\in Z^2(X_0,L_0)\).
采用的符号约定为 \(\mu'=\delta(\mu,\mu')\cdot\mu\);
即对每个 \(S'\to S\) 及 \(x,y\in X(S')\),
\[
 \mu'(x,y)=\delta(\mu,\mu')(x_0,y_0)\cdot\mu(x,y).
\]
\footnote{编者注: 我们沿用原文的符号约定,
以便在 \SGARefLink{sga3:iii:paragraph:4-38:item:5}{4.38(5)} 中有等式
\(\partial^1\bar d(X,X')=\bar\delta(X,X')\); 另见编者注 54.}

% Source page: Exp3 local p. 41; combined p. 151; original printed p. 126.
以 \(\bar\delta(\mu,\mu')\) 表示 \(\delta(\mu,\mu')\)
在 \(H^2(X_0,L_0)\) 中的像.
最后, 若将配备 \(\mu\) (分别地, \(\mu'\)) 的 \(X\)
简记为 \(X\) (分别地, \(X'\)), 则用 \(\delta(X,X')\)
代替 \(\delta(\mu,\mu')\), 对 \(\bar\delta(X,X')\) 也作类似约定.

\medskip
\noindent\SGAInlineRefTarget{sga3:iii:remark:3-6}\textbf{注 3.6.}
设 \(X_{\mathscr J}\) 是一个平滑仿射 \(S_{\mathscr J}\)-概形.
由 \SGARefLink{sga3:iii:proposition:0-15}{0.15}, 在同构意义下,
存在唯一的在 \(S\) 上平滑的 \(S\)-概形 \(X\),
其模 \(\mathscr J\) 的约化为 \(X_{\mathscr J}\).
若 \(X_{\mathscr J}\) 配备了一个 \(S_{\mathscr J}\)-群结构,
则由 \SGARefLink{sga3:iii:proposition:0-16}{0.16}, 条件 \((i_1)\) 自动成立.
此外, \SGARefLink{sga3:iii:remark:0-6}{0.6} 简化了 \(L_0\) 的定义, 并给出:

\medskip
\noindent\SGAInlineRefTarget{sga3:iii:corollary:3-7}\textbf{推论 3.7.}
\textit{设 \(S,\mathscr I,\mathscr J\) 如 \SGARefLink{sga3:iii:lemma:3-1}{3.1}
中所述, \(X_{\mathscr J}\) 是一个平滑仿射 \(S_{\mathscr J}\)-群.}
\begin{enumerate}
\item[(i)] \textit{模 \(\mathscr J\) 的约化为 \(X_{\mathscr J}\)
的平滑 \(S\)-群, 模去在 \(X_{\mathscr J}\) 上诱导恒等映射的同构后,
其集合或者为空, 或者是下列群的主齐性集合:}
\[
 H^2\!\left(
  X_0,\operatorname{Lie}(X_0/S_0)
  \otimes_{\mathscr O_{S_0}}\mathscr J
 \right).
\]
\item[(ii)] \textit{存在一个模 \(\mathscr J\) 的约化为
\(X_{\mathscr J}\) 的平滑 \(S\)-群,
当且仅当下列群中的某个障碍}
\[
 H^3\!\left(
  X_0,\operatorname{Lie}(X_0/S_0)
  \otimes_{\mathscr O_{S_0}}\mathscr J
 \right)
\]
\textit{为零.}
\end{enumerate}

照常, 得到以下推论.

\medskip
\noindent\SGAInlineRefTarget{sga3:iii:corollary:3-8}\textbf{推论 3.8.}
\textit{设 \(S\) 是一个概形, \(S_0\) 是由幂零理想 \(\mathscr I\)
定义的闭子概形. 设 \(X_0\) 是一个平滑仿射 \(S_0\)-群.}
\begin{enumerate}
\item[(i)] \textit{若对每个 \(n\geq0\), 有}
\[
 H^2\!\left(
  X_0,\operatorname{Lie}(X_0/S_0)
  \otimes_{\mathscr O_{S_0}}
  \mathscr I^{n+1}/\mathscr I^{n+2}
 \right)=0,
\]
\textit{则约化为 \(X_0\) 的任意两个平滑 \(S\)-群,
通过一个在 \(X_0\) 上诱导恒等映射的同构而同构.}
\item[(ii)] \textit{若对每个 \(n\geq0\), 有}
\[
 H^3\!\left(
  X_0,\operatorname{Lie}(X_0/S_0)
  \otimes_{\mathscr O_{S_0}}
  \mathscr I^{n+1}/\mathscr I^{n+2}
 \right)=0,
\]
\textit{则存在一个约化为 \(X_0\) 的平滑 \(S\)-群.}
\end{enumerate}

\medskip
\noindent\SGAInlineRefTarget{sga3:iii:corollary:3-9}\textbf{推论 3.9.}
\textit{设 \(S\) 仿射, \(S_0\) 是由幂零理想 \(\mathscr I\)
定义的闭子概形. 假设每个 \(\mathscr I^{n+1}/\mathscr I^{n+2}\)
在 \(S_0\) 上局部自由. 设 \(X_0\) 是一个平滑仿射 \(S_0\)-群.}
\begin{enumerate}
\item[(i)] \textit{若 \(H^2(X_0,\operatorname{Lie}(X_0/S_0))=0\),
则约化为 \(X_0\) 的任意两个平滑 \(S\)-群都同构.}
\item[(ii)] \textit{若 \(H^3(X_0,\operatorname{Lie}(X_0/S_0))=0\),
则存在一个约化为 \(X_0\) 的平滑 \(S\)-群.}
\end{enumerate}

% Source page: Exp3 local p. 42; combined p. 152; original printed p. 126.
\medskip
\noindent\SGAInlineRefTarget{sga3:iii:corollary:3-10}\textbf{推论 3.10.}
\textit{设 \(S_0\) 是一个概形, \(S=I_{S_0}\) 是
\(S_0\) 上的对偶数概形. 设 \(X_0\) 是一个平滑 \(S_0\)-群.
则对于每个满足 \(Y_0\) 与 \(X_0\) 在 \(S_0\) 上同构的平滑
\(S\)-群 \(Y\), 它都与 \(X=X_0\times_{S_0}S\) 在 \(S\) 上同构,
当且仅当}
\[
 H^2(X_0,\operatorname{Lie}(X_0/S_0))=0.
\]
\footnote{编者注: 此结果用于 XXIV, \SGARefLink{sga3:xxiv:1.13}{1.13}.}

事实上, 由 \SGARefLink{sga3:iii:theorem:3-5}{3.5},
将这样的群 \(Y\) 模去在 \(X_0\) 上诱导恒等映射的
\(S\)-群同构所得的类的集合,
与 \(H^2(X_0,\operatorname{Lie}(X_0/S_0))\) 之间有双射.
因此, 模去任意 \(S\)-群同构所得的类, 与
\[
 H^2(X_0,\operatorname{Lie}(X_0/S_0))/\Gamma_0,
\]
之间有双射, 其中
\[
 \Gamma_0=\operatorname{Aut}_{S_0\text{-}\mathrm{gr.}}(X_0),
\]
它以显然的方式作用于 \(H^2\). 立即可得结论.
\footnote{编者注: 事实上,
\(\operatorname{Aut}_{S_0\text{-}\mathrm{gr.}}(X_0)\)
通过群自同构作用于阿贝尔群 \(H^2(X_0,\operatorname{Lie}(X_0/S_0))\),
因此 \(0\) 的轨道是单元素集 \(\{0\}\).
所以, 商集是单元素集, 当且仅当
\(H^2(X_0,\operatorname{Lie}(X_0/S_0))=\{0\}\).}

% Continue with section 4. Its heading begins at the bottom of Exp3 local
% p. 42 (combined p. 152; original printed p. 127).

% Translation coverage: all of Section 3, through Corollary 3.10.
% Statistics: source lines=405; target lines=394; inline math=316; display math=13; equation/align/multline blocks=1; unresolved translation TODOs=0.
